\documentclass[11pt,a4paper]{article}
\usepackage[margin=2.5cm]{geometry}
\usepackage[T1]{fontenc}
\usepackage{booktabs}
\usepackage{xcolor}
\usepackage{amsmath}
\usepackage{listings}
\usepackage[hidelinks]{hyperref}

\emergencystretch=2em
\definecolor{codebg}{gray}{0.95}
\lstset{language=C++, basicstyle=\ttfamily\small, keywordstyle=\color{blue!70!black},
  commentstyle=\color{green!50!black}\itshape, backgroundcolor=\color{codebg},
  frame=single, rulecolor=\color{black!20}, breaklines=true, showstringspaces=false,
  tabsize=2, columns=flexible}

\title{ATmega32A Driver Course \\ \Large Module 5 --- ADC (Analog-to-Digital Converter)}
\author{Ali Sahafi}
\date{}

\begin{document}
\maketitle

\section{Theory}

The ADC converts a voltage into a number. The ATmega32A has one
\textbf{10-bit} ADC with an 8-channel multiplexer on \textbf{PORTA}
(ADC0--ADC7 = PA0--PA7): one converter, eight selectable inputs.

With 10 bits the result is 0--1023, measured relative to a
\textbf{reference voltage} $V_{\text{ref}}$:

\begin{equation*}
\text{result} = \frac{V_{\text{in}}}{V_{\text{ref}}} \cdot 1024
\qquad\Longleftrightarrow\qquad
V_{\text{in}} = \frac{\text{result}}{1024} \cdot V_{\text{ref}}
\end{equation*}

Three reference options: the \texttt{AREF} pin (\texttt{ADC\_REF\_AREF}),
the supply voltage AVCC (\texttt{ADC\_REF\_AVCC}, the usual choice) or the
internal 2.56\,V bandgap (\texttt{ADC\_REF\_INTERNAL}).

\paragraph{Prescaler.} The ADC needs a clock between 50 and 200\,kHz for full
resolution. The prescaler divides $F_{\text{CPU}}$ down into that window: at
8\,MHz, prescaler 64 gives 125\,kHz --- the driver default. One conversion
takes about 13 ADC clock cycles ($\approx$\,100\,$\mu$s).

\section{Registers}

\begin{center}
\small
\begin{tabular}{lll}
\toprule
\textbf{Register} & \textbf{Role} & \textbf{Bits used by the driver} \\
\midrule
\texttt{ADMUX}  & Reference + channel select & \texttt{REFS1:0}, \texttt{MUX2:0} \\
\texttt{ADCSRA} & Control and status & \texttt{ADEN}, \texttt{ADSC}, \texttt{ADATE}, \texttt{ADIE}, \texttt{ADPS2:0} \\
\texttt{ADCL/ADCH} & 10-bit result & read \texttt{ADCL} first, then \texttt{ADCH}! \\
\texttt{SFIOR}  & Auto-trigger source & \texttt{ADTS2:0} = 000 (free running) \\
\bottomrule
\end{tabular}
\end{center}

\section{API reference}

\begin{center}
\small
\begin{tabular}{ll}
\toprule
\textbf{Method} & \textbf{Description} \\
\midrule
\texttt{ADC.init(ref, prescaler, interrupt)} & Configure reference, prescaler, optional interrupt \\
\texttt{ADC.read(channel)} & Blocking read --- waits for conversion, returns 0--1023 \\
\texttt{ADC.startConversion(channel)} & Non-blocking single-shot conversion \\
\texttt{ADC.startContinuous(channel)} & Free-running mode --- auto-converts continuously \\
\texttt{ADC.stopContinuous()} & Stop free-running mode \\
\texttt{ADC.isReady()} & Returns \texttt{true} when conversion is complete \\
\texttt{ADC.lastResult()} & Read last result without triggering a new conversion \\
\texttt{ADC.stop()} / \texttt{ADC.start()} & Disable / re-enable the ADC \\
\texttt{ADC.attachInterrupt(callback)} & Set conversion callback: \texttt{void cb(uint16\_t value)} \\
\bottomrule
\end{tabular}
\end{center}

\begin{itemize}
  \item \textbf{Reference:} \texttt{ADC\_REF\_AREF}, \texttt{ADC\_REF\_AVCC}, \texttt{ADC\_REF\_INTERNAL}
  \item \textbf{Prescaler:} \texttt{ADC\_PRESCALER\_2/4/8/16/32/64/128}
  \item \textbf{Channel:} \texttt{ADC\_CHANNEL\_0} \dots\ \texttt{ADC\_CHANNEL\_7}
\end{itemize}

Three usage styles, from simplest to most advanced:
\begin{itemize}
  \item \textbf{blocking:} \texttt{read()} --- the CPU waits for the conversion
  \item \textbf{polling:} \texttt{startConversion()}, then \texttt{isReady()} + \texttt{lastResult()}
  \item \textbf{interrupt-driven:} \texttt{startContinuous()} + callback --- see the
        capstone \texttt{example.cpp} in the repository root
\end{itemize}

\section{Example: potentiometer bar graph}

\paragraph{Wiring.}
\begin{itemize}
  \item Potentiometer ($\geq$1\,k$\Omega$): outer pins to VCC and GND, wiper to \textbf{PA0}.
  \item 8 LEDs + 330\,$\Omega$ resistors from \textbf{PC0--PC7} to GND.
  \item \textbf{AVCC must be connected to VCC} --- the ADC is powered from the AVCC pin.
\end{itemize}

\paragraph{Build \& flash.} \texttt{make adc}

\lstinputlisting{example.cpp}

\section{Exercises}

\begin{enumerate}
  \item The pot is at mid-position. What value does \texttt{ADC.read()}
        return (roughly), and what voltage does that correspond to at
        $V_{\text{ref}} = 5$\,V?
  \item Rewrite the example in the \emph{polling} style
        (\texttt{startConversion} / \texttt{isReady} / \texttt{lastResult})
        so the LED bar updates while the main loop stays free to blink a
        ninth ``heartbeat'' LED.
  \item Connect an LM35 temperature sensor (10\,mV/$^{\circ}$C) to ADC1 and
        compute the temperature in degrees. Which reference gives better
        resolution for this sensor: AVCC (5\,V) or internal 2.56\,V?
\end{enumerate}

\end{document}
